% Generated draft: requires mathematical and source review.
\documentclass[11pt]{article}
\usepackage[margin=25mm]{geometry}
\usepackage{fontspec}
\setmainfont{Latin Modern Roman}
\usepackage{amsmath,amssymb,mathtools,unicode-math}
\setmathfont{Latin Modern Math}
\usepackage{enumitem,needspace,xurl,hyperref}
\setlength{\parindent}{0pt}
\setlength{\parskip}{5pt}
\setlength{\emergencystretch}{3em}
\newcommand{\N}{\mathbb{N}}
\newcommand{\Z}{\mathbb{Z}}
\newcommand{\Q}{\mathbb{Q}}
\newcommand{\R}{\mathbb{R}}
\newcommand{\C}{\mathbb{C}}
\newcommand{\F}{\mathbb{F}}
\newcommand{\A}{\mathbb{A}}
\newcommand{\PP}{\mathbb{P}}
\newcommand{\E}{\mathbb{E}}
\newcommand{\eps}{\varepsilon}
\newcommand{\dd}{\,\mathrm{d}}
\begin{document}
\section*{UnsolvedMath ID: 30001149}
\textbf{Source problem code:} catalogue entry
\subsection{Definitions and mathematical statement}
For a field \(k\), let \(k[[t]]\) denote the ring of formal power series in one variable over \(k\). Let \(p\) be a prime and \(n\geq 2\) an integer, so that \(p^n>p\). An automorphism \(\sigma\) has exact order \(p^n\) if
\[
\sigma^{p^n}=\operatorname{id},\qquad
\sigma^j\neq\operatorname{id}\quad(1\leq j<p^n).
\]
The supplied record does not specify the hypotheses on \(k\), including its characteristic, or whether the automorphisms must fix \(k\) and be continuous for the \(t\)-adic topology. It also does not define the source-specific property ``almost rational.'' Consequently, the following is a formulation schema rather than a fully specified classification problem.

Write \(G_k\) for the automorphism group intended by the source and \(\mathcal A_k(\sigma)\) for its intended almost-rationality condition; these symbols are placeholders for definitions absent from the record. The requested problem is to classify
\[
\mathcal C_{k,p,n}
=\left\{\sigma\in G_k:
\mathcal A_k(\sigma),\quad
\sigma^{p^n}=\operatorname{id},\quad
\sigma^j\neq\operatorname{id}\ \text{for every }1\leq j<p^n
\right\}
\]
for the fields and primes allowed by the original source, without imposing the additional condition that every power \(\sigma^j\), for \(j\neq 1\), is almost rational with respect to the same Artin--Schreier root.

For terminology only, over a field \(F\) of characteristic \(p\), an Artin--Schreier root is an element \(\beta\) in an extension field satisfying \(\beta^p-\beta=a\) for some \(a\in F\). The supplied record does not identify the relevant field \(F\), equation, or relation between such a root and almost rationality. It also leaves unspecified the range of the exponent \(j\) in the omitted additional condition and the equivalence relation, if any, under which the classification is requested. These qualifications must be resolved from the source before a fully precise statement can be given.
\subsection{Short English statement}
Classify almost-rational automorphisms of \(k[[t]]\) having exact order \(p^n\), with \(p\) prime and \(n\geq2\), without requiring their other powers to be almost rational with respect to one common Artin--Schreier root. The supplied record omits definitions and field hypotheses needed to make this classification target fully precise.
\subsection{Sources}
\begin{itemize}
\item Local UnsolvedMath registry, supplied complete source record: permanent ID \(30001149\), source problem code \texttt{catalogue entry}, title ``Almost-Rational Automorphisms of Higher Prime-Power Order,'' in the Oberwolfach Reports Open Problems collection. The record gives proposed year 2009 and names no proposer.
\item Supplied source citation: \emph{The Arithmetic of Fields} (2009), Oberwolfach Reports, pp.~285--354, DOI \texttt{10.4171/owr/2009/05}. \url{https://doi.org/10.4171/owr/2009/05}. The record describes this report as documentation of a mathematical workshop and its open-problem discussions.
\item Supplied literature source labeled ``Bleher--Poonen--Pop--Symonds'': \url{https://oa.tib.eu/renate/bitstreams/8fd427cc-2386-49b6-bb25-b696af831780/download}. No title or further bibliographic details are supplied for this reference.
\item Literature assessment dated exactly \texttt{2026-08-21}: the supplied curation labels the problem open, states that the source leaves the classification open without conditions on the other powers, and reports that no later complete classification was found. This is supplied dated curation, not an independently verified assessment. The record labels its statement \texttt{corrected\_verified} and describes a meaning-preserving semantic, typographic, and LaTeX cleanup; those labels do not establish mathematical correctness or current literature status.
\end{itemize}
\end{document}
